\section{Conclusion}
\label{sec:conclusion}

Design fingerprinting via within-game z-score normalization provides a quality-neutral
representation of a game's characteristic design emphasis. The normalization reduces PC1
from 81.3\% to 50.0\%, surpassing the pre-registered threshold of $\leq 55\%$, and
produces mean pairwise Euclidean distances of 3.9 (range 1.1--8.7) between fingerprints,
confirming that games genuinely differ in design character across the 32-axis space even
when they occupy similar quality tiers.

The fingerprint is both an analytical tool (recovering known design distinctions from
raw scoring data) and a design diagnostic (identifying which axes a design is
deprioritizing relative to its own average). The method is general and can be applied
to any multi-attribute evaluation dataset where a general quality factor dominates
raw PCA.

All fingerprints for 150 corpus games, analysis scripts, and visualization code are
available at [URL].
